\documentclass[12pt,letterpaper]{article}

\usepackage[margin=0.75in]{geometry}
\usepackage{amsmath,amssymb}
\usepackage{enumitem}

\pagestyle{empty}
\setlength{\parindent}{0pt}
\setlength{\parskip}{0.5em}
\setlength{\emergencystretch}{2em}

\begin{document}

\begin{center}
  {\large\bfseries MATH 15200: Calculus II}\\
  {\large\bfseries Homework 4: Logarithms and Exponentials}
\end{center}

\vspace{0.25em}
\hrule
\vspace{0.8em}

\textbf{Due:} Monday, November 9, 2026, at 11:59 PM.

This is the written portion of Homework 4; complete the routine online portion
separately in Edfinity. Upload your written solutions to Gradescope, via Canvas. You may
collaborate, but you must write and submit your own solutions independently.



\vspace{0.8em}


\begin{enumerate}[label=\textbf{\arabic*.},leftmargin=2em,itemsep=1.25em]
  \item For each positive integer $n$, let \(H_n\) be the harmonic sum
    \[
      H_n=1+\frac12+\frac13+\cdots+\frac1n.
    \]
    Show that \(H_n\) grows logarithmically by proving that 
    \[
      \ln(n+1)\le H_n\le 1+\ln n.
    \]

    \emph{Hint:} If $k\le t\le k+1$, then $\frac{1}{k+1}\le\frac1t\le\frac1k$.
    Use Problem~2 of Homework~2 to integrate these inequalities over
    $[k,k+1]$. Then add the results: over $k=1,\dots,n$ for the lower bound,
    and over $k=1,\dots,n-1$ for the upper bound.

  \item Without using a calculator, determine which is larger, $e^\pi$ or
    $\pi^e$.

    \emph{Hint:} Since $\ln$ is increasing, it is enough to compare
    $\ln\bigl(e^\pi\bigr)$ and $\ln\bigl(\pi^e\bigr)$. Rewrite this as a
    comparison of $g(e)$ and $g(\pi)$, where $g(x)=\frac{\ln x}{x}$ for
    $x>0$. Find where $g$ is increasing and where it is decreasing.

  \item Let $b>0$, and let $R$ be the region bounded by the curve
    $y=e^{-x^2}$, the $x$-axis, the $y$-axis, and the line $x=b$. The function
    $e^{-x^2}$ has no antiderivative given by a formula in familiar functions,
    so the area of $R$ cannot be computed directly with the Fundamental
    Theorem of Calculus.
    \begin{enumerate}[label=(\alph*)]
      \item Find the volume $V(b)$ of the solid formed by revolving $R$ about
        the $y$-axis.
      \item Find $\displaystyle\lim_{b\to\infty}V(b)$.
    \end{enumerate}

    \emph{Hint:} Use the shell method, and then the substitution $u=x^2$.
\end{enumerate}

\end{document}
